% ============================================================================
% CAPA INSTITUCIONAL — EMPREENDEDORISMO TECNOLÓGICO E DESENVOLVIMENTO TERRITORIAL
% ============================================================================
\begin{titlepage}
\thispagestyle{empty}
\begin{tikzpicture}[remember picture, overlay]
  \fill[etIndigoDark] (current page.north west) rectangle ([yshift=-16.5cm]current page.north east);
  \fill[etTeal] ([yshift=-16.5cm]current page.north west) rectangle ([yshift=-16.95cm]current page.north east);
  \fill[etAmber] ([yshift=-16.95cm]current page.north west) rectangle ([yshift=-17.2cm]current page.north east);
  \fill[bgSoft] ([yshift=-17.2cm]current page.north west) rectangle (current page.south east);
  \fill[etTeal, opacity=0.22] ([xshift=16.5cm, yshift=-3.2cm]current page.north west) circle (6.0cm);
  \fill[etAmber, opacity=0.14] ([xshift=19.0cm, yshift=-11.8cm]current page.north west) circle (4.6cm);
  \draw[white, opacity=0.16, line width=1pt] ([xshift=1.8cm, yshift=-1.8cm]current page.north west) rectangle ([xshift=-1.8cm, yshift=1.8cm]current page.south east);
\end{tikzpicture}

\vspace*{1.2cm}
\begin{center}
  {\color{etAmber}\bfseries\large\sffamily UNIVERSIDADE ESTADUAL DE FEIRA DE SANTANA (UEFS)}\\[0.25cm]
  {\color{bgAmber}\normalsize\sffamily Programa de Pós-Graduação em Planejamento Territorial (PPGPT)}\\[1.4cm]

  {\color{white}\bfseries\fontsize{28}{34}\selectfont\sffamily EMPREENDEDORISMO TECNOLÓGICO}\\[0.35cm]
  {\color{etAmber}\bfseries\fontsize{22}{28}\selectfont\sffamily E DESENVOLVIMENTO TERRITORIAL}\\[1.0cm]

  \begin{tcolorbox}[colback=white!10!etIndigoDark, colframe=etAmber, arc=4pt, boxrule=1pt, width=0.90\textwidth, halign=center, top=10pt, bottom=10pt]
    {\color{white}\large\sffamily Ecossistemas de Inovação, Impacto Territorial, Modelos de Negócio,\\
    Decisão Locacional, Análise de Investimentos e \textit{Valuation}}
  \end{tcolorbox}
  \vspace{1.0cm}

  {\color{bgAmber}\normalsize Apostila completa da disciplina (60 h $\bullet$ 15 encontros) $\bullet$ Explicação passo a passo sem pular contas\\
  Estudos de caso territoriais e quantitativos resolvidos $\bullet$ Gabarito integralmente comentado\\
  Simulados por unidade e roteiro completo para AV1, AV2 e Projeto Final (AV3)}
\end{center}

\vspace{3.0cm}

\begin{center}
  \begin{tcolorbox}[colback=white, colframe=etIndigo, arc=4pt, boxrule=1pt, width=0.90\textwidth, left=14pt, right=14pt, top=12pt, bottom=12pt]
    \begin{tabularx}{\textwidth}{@{}l >{\raggedright\arraybackslash}X@{}}
      \textbf{\color{etIndigo}Programa / Instituição:} & PPGPT — Universidade Estadual de Feira de Santana (UEFS) \\
      \textbf{\color{etIndigo}Docente Responsável:} & Prof. Dr. Luiz Diego Vidal Santos (`ldvsantos@uefs.br`) \\
      \textbf{\color{etIndigo}Carga Horária:} & 60 horas-aula (15 encontros de 4 horas) \\
      \textbf{\color{etIndigo}Estrutura Didática:} & 4 Unidades $\bullet$ 15 Capítulos em 10 Etapas $\bullet$ 4 Simulados $\bullet$ Formulário e Glossário \\
    \end{tabularx}
  \end{tcolorbox}
\end{center}

\vfill
\begin{center}
  {\small\color{etMuted}\bfseries FEIRA DE SANTANA — BAHIA $\bullet$ EDIÇÃO A4 PRONTA PARA IMPRESSÃO E ESTUDO AUTÔNOMO}
\end{center}
\end{titlepage}

% ============================================================================
% APRESENTAÇÃO DA DISCIPLINA E GUIA DE USO
% ============================================================================
\cleardoublepage
\phantomsection
\addcontentsline{toc}{chapter}{Apresentação da Disciplina e Guia de Uso}
\chapter*{Apresentação da Disciplina e Guia de Uso da Apostila}
\markboth{APRESENTAÇÃO E GUIA DE USO}{APRESENTAÇÃO E GUIA DE USO}

Bem-vindo(a) à apostila completa da disciplina \textbf{Empreendedorismo Tecnológico e Desenvolvimento Territorial}, ofertada no âmbito do Programa de Pós-Graduação em Planejamento Territorial (PPGPT) da Universidade Estadual de Feira de Santana (UEFS).

Este livro didático foi projetado para articular rigor teórico-crítico sobre território e inovação com domínio operacional de modelagem de negócios, avaliação locacional e engenharia econômico-financeira de projetos. Ele parte dos conceitos fundadores até a resolução passo a passo de fluxos de caixa, VPL, TIR, *Valuation* e construção do Projeto Final Integrador (AV3).

\section*{Arquitetura da Disciplina (60 h $\bullet$ 15 Encontros)}

\begin{center}
\renewcommand{\arraystretch}{1.24}
\begin{tabularx}{\textwidth}{@{}l >{\raggedright\arraybackslash}X c c@{}}
  \toprule
  \rowcolor{etIndigo}
  \textcolor{white}{\bfseries Unidade} & \textcolor{white}{\bfseries Eixo Temático e Capítulos} & \textcolor{white}{\bfseries Carga} & \textcolor{white}{\bfseries Avaliação} \\
  \midrule
  \textbf{Unidade I} & Território e Ecossistemas (Caps. 1 a 6): EBT, barreiras, APLs, SRI, Hélice, Empreendedorismo Social, Tecnologia Social, Bioeconomia e IGs. & 24 h & \textbf{AV1 (30\%)} \\
  \rowcolor{bgTeal}
  \textbf{Unidade II} & Processo e Estratégia (Caps. 7 a 11): Effectuation, Vale da Morte, TAM/SAM/SOM, PI, Canvas, Planos, Fomento Regional e Spin-offs. & 20 h & Base da AV3 \\
  \textbf{Unidade III} & Projetos e Viabilidade (Caps. 12 a 14): Decisão Locacional, CAPEX/OPEX, VPL, TIR, MTIR, VAUE, Payback e Valuation (FCD/VC). & 12 h & \textbf{AV2 (20\%)} \\
  \rowcolor{bgAmber}
  \textbf{Unidade IV} & Síntese Territorial (Cap. 15): Seminário Final Integrador e Defesa Oral (20 min + 10 min de arguição). & 4 h & \textbf{AV3 (50\%)} \\
  \bottomrule
\end{tabularx}
\end{center}

\section*{Sistema Oficial de Avaliação}

\begin{itemize}[leftmargin=1.5em, itemsep=3pt]
  \item \textbf{AV1 (30\%):} Seminário e fichamento crítico sobre os debates e contrapontos da Unidade I (Caps. 1 a 6).
  \item \textbf{AV2 (20\%):} Estudo dirigido quantitativo e locacional sobre elaboração de projetos, decisão locacional, VPL, TIR, Payback e *Valuation* (Caps. 12 a 14).
  \item \textbf{AV3 (50\%):} Trabalho final aplicado + defesa oral (20 min) — Diagnóstico de um ecossistema territorial ou proposta de empreendimento tecnológico/de impacto com viabilidade locacional e econômico-financeira (rubrica: 30\% Leitura Territorial, 30\% Estratégia/Modelo, 25\% Viabilidade/\textit{Valuation}, 15\% Comunicação/Defesa).
\end{itemize}

\section*{Método Pedagógico em 10 Etapas}

Cada um dos 15 capítulos segue rigorosamente a mesma estrutura didática: (1) Objetivos; (2) Revisão de pré-requisitos; (3) Explicação teórica com vocabulário prévio, figuras e dedução de modelos/fórmulas; (4) 5 Estudos de Caso/Exemplos resolvidos em 7 passos sem omitir contas; (5) 4 Erros Mais Comuns; (6) 15 Exercícios em 3 níveis com espaço impresso; (7) 3 Questões estilo ANPAD/POSCOMP/concursos de CT\&I e desenvolvimento regional; (8) Gabarito integralmente comentado; (9) Resumo do capítulo; (10) Teste de aprendizagem (Checkpoint).
